% !TeX program = pdflatex
\documentclass[11pt,a4paper]{amsart}

\newcommand{\notelanguage}{finnish}
% Shared preamble for course notes and exercise sets.

\usepackage[T1]{fontenc}
\usepackage{lmodern}
\providecommand{\notelanguage}{english}
\usepackage[finnish,english]{babel}
\AtBeginDocument{\expandafter\selectlanguage\expandafter{\notelanguage}}

\newcommand{\notesfinnishlanguage}{finnish}
\ifx\notelanguage\notesfinnishlanguage
\newcommand{\notestheoremname}{Lause}
\newcommand{\noteslemmaname}{Lemma}
\newcommand{\notespropositionname}{Propositio}
\newcommand{\notescorollaryname}{Seurauslause}
\newcommand{\notesdefinitionname}{Määritelmä}
\newcommand{\notesexamplename}{Esimerkki}
\newcommand{\notesproblemname}{Tehtävä}
\newcommand{\notesexercisename}{Tehtävä}
\newcommand{\notesremarkname}{Huomautus}
\newcommand{\notessolutionname}{Ratkaisu}
\else
\newcommand{\notestheoremname}{Theorem}
\newcommand{\noteslemmaname}{Lemma}
\newcommand{\notespropositionname}{Proposition}
\newcommand{\notescorollaryname}{Corollary}
\newcommand{\notesdefinitionname}{Definition}
\newcommand{\notesexamplename}{Example}
\newcommand{\notesproblemname}{Problem}
\newcommand{\notesexercisename}{Exercise}
\newcommand{\notesremarkname}{Remark}
\newcommand{\notessolutionname}{Solution}
\fi

\usepackage[a4paper,margin=30mm]{geometry}
\usepackage{microtype}
\usepackage{graphicx}
\usepackage{enumitem}

\newcommand{\exerciseheading}[2]{%
  \par\noindent
  {\large\bfseries #1\par}
  \vspace{0.1em}
  \noindent{\large #2\par}
  \vspace{1.5em}
}

\usepackage{amsmath}
\usepackage{amssymb}
\usepackage{amsthm}
\usepackage{mathtools}
\usepackage{aliascnt}

% Linear systems: the argument gives the number of alignment pairs.
\newenvironment{linearsystem}[1]
  {\left\{\begin{alignedat}{#1}}
  {\end{alignedat}\right.}

\usepackage[hidelinks,pdfusetitle]{hyperref}

\numberwithin{equation}{section}
\setlist{itemsep=0.25em,topsep=0.5em}

\theoremstyle{plain}
\newtheorem{theorem}{\notestheoremname}[section]
\newaliascnt{lemma}{theorem}
\newtheorem{lemma}[lemma]{\noteslemmaname}
\aliascntresetthe{lemma}
\newaliascnt{proposition}{theorem}
\newtheorem{proposition}[proposition]{\notespropositionname}
\aliascntresetthe{proposition}
\newaliascnt{corollary}{theorem}
\newtheorem{corollary}[corollary]{\notescorollaryname}
\aliascntresetthe{corollary}

\theoremstyle{definition}
\newaliascnt{definition}{theorem}
\newtheorem{definition}[definition]{\notesdefinitionname}
\aliascntresetthe{definition}
\newaliascnt{example}{theorem}
\newtheorem{example}[example]{\notesexamplename}
\aliascntresetthe{example}
\newaliascnt{problem}{theorem}
\newtheorem{problem}[problem]{\notesproblemname}
\aliascntresetthe{problem}
\newtheorem{exercise}{\notesexercisename}

\theoremstyle{remark}
\newaliascnt{remark}{theorem}
\newtheorem{remark}[remark]{\notesremarkname}
\aliascntresetthe{remark}

\newenvironment{solution}
{
\begin{proof}[\notessolutionname]}
  {
\end{proof}}

\usepackage[nameinlink,noabbrev]{cleveref}

\crefname{theorem}{theorem}{theorems}
\crefname{lemma}{lemma}{lemmas}
\crefname{proposition}{proposition}{propositions}
\crefname{corollary}{corollary}{corollaries}
\crefname{definition}{definition}{definitions}
\crefname{example}{example}{examples}
\crefname{problem}{problem}{problems}
\crefname{exercise}{exercise}{exercises}
\crefname{remark}{remark}{remarks}

\ifx\notelanguage\notesfinnishlanguage
\crefname{theorem}{lause}{lauseet}
\crefname{lemma}{lemma}{lemmat}
\crefname{proposition}{propositio}{propositiot}
\crefname{corollary}{seurauslause}{seurauslauseet}
\crefname{definition}{määritelmä}{määritelmät}
\crefname{example}{esimerkki}{esimerkit}
\crefname{problem}{tehtävä}{tehtävät}
\crefname{exercise}{tehtävä}{tehtävät}
\crefname{remark}{huomautus}{huomautukset}
\fi

\newcommand{\NN}{\mathbb{N}}
\newcommand{\NNone}{\mathbb{N}_1}
\newcommand{\ZZ}{\mathbb{Z}}
\newcommand{\QQ}{\mathbb{Q}}
\newcommand{\RR}{\mathbb{R}}
\newcommand{\CC}{\mathbb{C}}
\newcommand{\PP}{\mathbb{P}}

\DeclareMathOperator{\im}{im}
\DeclareMathOperator{\rank}{rank}
\DeclareMathOperator{\syt}{syt}
\DeclarePairedDelimiter{\abs}{\lvert}{\rvert}
\DeclarePairedDelimiter{\norm}{\lVert}{\rVert}
\newcommand{\set}[1]{\left\lbrace#1\right\rbrace}
\newcommand{\maxset}[1]{\max\set{#1}}
\newcommand{\minset}[1]{\min\set{#1}}
\newcommand{\powerset}[1]{\mathcal{P}\!\left(#1\right)}


\usepackage{tikz}

\title{Analyysi II\\Ryhmätyö 1 ratkaisut}
\author{}
\date{}

\begin{document}

\exerciseheading{Analyysi II}{Ryhmätyö 1 ratkaisut}

\begin{exercise}
  Tutki, missä pisteissä funktio \(f\colon [0,3] \to \RR\),
  \[
    f(x) =
    \begin{cases}
      \sin x, & 0 < x < \frac{\pi}{2},\\
      1, & x = 0 \text{ tai } \frac{\pi}{2} \leqslant x < 2,\\
      x - 2, & 2 \leqslant x \leqslant 3,
    \end{cases}
  \]
  \begin{enumerate}[label=(\alph*), nosep]
    \item on vasemmalta jatkuva,
    \item on oikealta jatkuva,
    \item on jatkuva,
    \item ei ole jatkuva.
  \end{enumerate}
  Piirrä funktion kuvaaja. Tässä tehtävässä kuvasta katsominen riittää
  perusteluksi! Saat tarkastella sinifunktiota \(x \mapsto \sin x\)
  lukiotietojen perusteella (sinifunktion määritelmä ja jatkuvuus esitetään
  oppikirjan luvuissa 1.6 ja 4.2; niitä käsitellään tällä kurssilla myöhemmin).

  Huom! Tee jatkuvuustarkastelut \emph{pisteittäin} Määritelmien 4.1.1 ja
  4.1.4 perusteella, älä käytä Määritelmässä 4.1.8 esitettyä jatkuvuutta
  \emph{joukossa}.
\end{exercise}

\begin{solution}
  \leavevmode\par
  \begin{enumerate}[label=(\alph*), nosep]
    \item Koska \(x \mapsto \sin x\) on jatkuva jokaisessa pisteessä
      \(x \in \RR\), myös \(f\) on jatkuva ja siten vasemmalta jatkuva
      jokaisessa pisteessä \(x \in (0, \pi/2)\).
      Lisäksi
      \[
        f(\pi/2) = 1 = \sin(\pi/2) = \lim_{x \to \pi/2^-} f(x),
      \]
      joten \(f\) on vasemmalta jatkuva myös pisteessä \(\pi/2\).
      Koska vakiofunktio \(x \mapsto 1\) on jatkuva jokaisessa pisteessä,
      \(f\) on vasemmalta jatkuva jokaisessa pisteessä
      \(x \in (\pi/2, 2)\). Samoin funktion \(x \mapsto x - 2\)
      jatkuvuudesta seuraa, että \(f\) on vasemmalta jatkuva jokaisessa
      pisteessä \(x \in (2, 3)\). Lisäksi \(f\) on vasemmalta jatkuva
      pisteessä \(3\), sillä
      \[
        f(3) = 1 = \lim_{x \to 3^-} f(x).
      \]
      Pisteessä \(2\) funktio ei ole vasemmalta jatkuva, sillä
      \(\lim_{x \to 2^-} f(x) = 1 \neq 0 = f(2)\).
      Pisteessä \(0\) vasemmalta jatkuvuutta ei ole määritelty, koska
      määrittelyjoukossa ei ole pisteitä nollan vasemmalla puolella.
      Siis \(f\) on vasemmalta jatkuva täsmälleen pisteissä
      \[
        x \in (0, 2) \cup (2, 3].
      \]
    \item Kuten kohdassa (a), sinifunktion jatkuvuudesta seuraa, että
      \(f\) on oikealta jatkuva jokaisessa pisteessä
      \(x \in (0, \pi/2)\).
      Lisäksi \(f\) on oikealta jatkuva pisteessä \(\pi/2\), sillä
      \[
        f(\pi/2) = 1 = \lim_{x \to (\pi/2)^+} f(x).
      \]
      Jokaisessa pisteessä \(x \in (\pi/2, 2)\) funktio \(f\) on
      oikealta jatkuva, koska se on tällä välillä vakio.
      Funktio \(f\) ei ole oikealta jatkuva pisteessä \(0\), sillä
      \[
        f(0) = 1 \neq 0 = \lim_{x \to 0^+} f(x).
      \]
      Funktion \(x \mapsto x - 2\) jatkuvuudesta seuraa, että \(f\) on
      oikealta jatkuva jokaisessa pisteessä \(x \in (2, 3)\).
      Lisäksi \(f\) on oikealta jatkuva pisteessä \(2\), sillä
      \[
        f(2) = 0 = \lim_{x \to 2^+} f(x).
      \]
      Pisteessä \(3\) oikealta jatkuvuutta ei ole määritelty, koska
      määrittelyjoukko ei sisällä pisteitä \(3\):n oikealla puolella.
      Siis \(f\) on oikealta jatkuva täsmälleen pisteissä
      \[
        x \in (0, 3).
      \]
    \item Funktio \(f\) on jatkuva sisäpisteessä täsmälleen silloin, kun
      se on sekä oikealta että vasemmalta jatkuva. Kohtien (a) ja (b)
      nojalla \(f\) on siis jatkuva täsmälleen pisteissä
      \[
        x \in (0, 2) \cup (2, 3).
      \]
    \item Kohdan (c) nojalla \(f\) ei ole jatkuva täsmälleen pisteissä
      \(x \in \set{0, 2, 3}\). Päätepisteissä \(0\) ja \(3\)
      kaksipuolista jatkuvuutta ei ole määritelty, ja pisteessä \(2\)
      toispuoleiset raja-arvot eroavat toisistaan.
  \end{enumerate}
  \begin{center}
    \begin{tikzpicture}[x=2.6cm,y=2.6cm,>=stealth]
      \draw[->] (-0.2,0) -- (3.35,0) node[right] {\(x\)};
      \draw[->] (0,-0.2) -- (0,1.4) node[above] {\(y\)};
      \draw[densely dashed,gray] ({pi/2},0) -- ({pi/2},1);
      \draw[densely dashed,gray] (2,0) -- (2,1);
      \draw[densely dashed,gray] (3,0) -- (3,1);
      \node[below left] at (0,0) {\(0\)};
      \foreach \x/\label in {1/1,1.5707963268/{\frac{\pi}{2}},2/2,3/3} {
        \draw (\x,0.025) -- (\x,-0.025) node[below] {\(\label\)};
      }
      \draw (0.025,1) -- (-0.025,1) node[left] {\(1\)};
      \draw[thick,domain=0:pi/2,samples=80] plot (\x,{sin(\x r)});
      \draw[thick] ({pi/2},1) -- (2,1);
      \draw[thick] (2,0) -- (3,1);
      \draw[thick,fill=white] (0,0) circle (2.5pt);
      \draw[thick,fill=white] (2,1) circle (2.5pt);
      \fill (0,1) circle (2.5pt);
      \fill ({pi/2},1) circle (2.5pt);
      \fill (2,0) circle (2.5pt);
      \fill (3,1) circle (2.5pt);
    \end{tikzpicture}
  \end{center}
\end{solution}

\begin{exercise}
  Olkoon \(f\colon \RR \to \RR\) sellainen funktio, että
  \[
    \lim_{h \to 0} [f(x + h) - f(x - h)] = 0
  \]
  kaikilla \(x \in \RR\). Onko \(f\) välttämättä jatkuva?
\end{exercise}

\end{document}
